\documentclass[12pt,a4paper]{article}

% ── Packages ──────────────────────────────────────────────────────────────────
\usepackage[margin=2.5cm, headheight=15pt]{geometry}
\usepackage{amsmath,amssymb,amsfonts}
\usepackage{graphicx}
\usepackage{booktabs}
\usepackage{array}
\usepackage{tabularx}
\usepackage{multirow}
\usepackage{xcolor}
\usepackage{hyperref}
\usepackage{enumitem}
\usepackage{titlesec}
\usepackage{fancyhdr}
\usepackage{natbib}
\usepackage{tikz}
\usetikzlibrary{shapes.geometric, arrows.meta, positioning, fit, backgrounds}
\usepackage{mdframed}
\usepackage{tcolorbox}
\tcbuselibrary{breakable,skins}
\usepackage{caption}
\usepackage{subcaption}

% ── Colours ───────────────────────────────────────────────────────────────────
\definecolor{accentgreen}{RGB}{0,128,96}
\definecolor{darkblue}{RGB}{0,51,102}
\definecolor{lightgray}{RGB}{245,245,245}
\definecolor{novelred}{RGB}{180,30,30}
\definecolor{existblue}{RGB}{30,80,160}

% ── Header / Footer ───────────────────────────────────────────────────────────
\pagestyle{fancy}
\fancyhf{}
\rhead{\textcolor{darkblue}{\small Network Anomaly Detection --- Research Proposal}}
\lhead{\textcolor{darkblue}{\small Tesfay Hagos}}
\rfoot{\thepage}
\renewcommand{\headrulewidth}{0.4pt}

% ── Section Formatting ────────────────────────────────────────────────────────
\titleformat{\section}{\large\bfseries\color{darkblue}}{\thesection.}{0.5em}{}[\titlerule]
\titleformat{\subsection}{\normalsize\bfseries\color{darkblue}}{\thesubsection.}{0.5em}{}

% ── Hyperref ──────────────────────────────────────────────────────────────────
\hypersetup{
    colorlinks=true,
    linkcolor=darkblue,
    citecolor=darkblue,
    urlcolor=accentgreen
}

% ─────────────────────────────────────────────────────────────────────────────
\begin{document}

% ── Title Page ────────────────────────────────────────────────────────────────
\begin{titlepage}
\centering
\vspace*{1cm}

{\color{darkblue}\rule{\linewidth}{2pt}}\\[0.5em]
{\LARGE\bfseries\color{darkblue}
Two-Pass Reconstruction-Error-Guided Attention\\[0.3em]
for Adversarial Autoencoder-Based\\[0.3em]
Network Anomaly Detection}\\[0.3em]
{\color{darkblue}\rule{\linewidth}{2pt}}

\vspace{1.5cm}

{\large\textbf{Research Proposal}}\\[0.5em]
{\normalsize Deep Learning Project --- Network Anomaly Detection}

\vspace{1.5cm}

{\large\textbf{Author:} Tesfay Hagos}\\[0.3em]
{\normalsize \href{mailto:tesfay.hagos1421@gmail.com}{tesfay.hagos1421@gmail.com}}

\vspace{0.8cm}
{\normalsize \textbf{Date:} June 2026}

\vspace{2cm}

\begin{abstract}
\noindent
Network intrusion detection requires identifying rare, unseen attack patterns in
high-dimensional traffic data with minimal or no labels.
Established unsupervised approaches on the KDD99 benchmark---deep autoencoders
\citep{ae_analysis_nids} and GAN-based discriminator scoring \citep{alad}---share
two structural limitations: all input features are treated uniformly regardless
of their discriminative relevance, and the latent representation is unconstrained,
making anomaly reasoning in latent space ill-defined.

This proposal describes an architectural experiment that addresses both limitations
through a single novel mechanism: \textbf{Two-Pass Reconstruction-Error-Guided Attention
(RE-Attention)}.
In the first pass, a standard encoder--decoder reconstructs the input; the
per-feature reconstruction error vector is then used as the attention signal to
re-weight the input for a second encoding pass whose latent code is regularised
by an adversarial Gaussian prior.
This creates a self-corrective feedback loop in which the model identifies its
own reconstruction failures and attends to them, enabling a richer and more
targeted anomaly representation than input-driven or bottleneck-driven attention
reported in prior work \citep{attention_ae_latent,vae_attention_lstm}.
A controlled five-condition ablation study isolates the contribution of each
component, and the attention weights are visualised per attack family to provide
interpretable insights into which network traffic features most strongly signal
each class of intrusion.
\end{abstract}

\vfill
{\small \textit{This proposal is for an experimental implementation study.
It is not submitted for publication.}}
\end{titlepage}

% ─────────────────────────────────────────────────────────────────────────────
\tableofcontents
\newpage

% =============================================================================
\section{Introduction}
% =============================================================================

\subsection{Background and Motivation}

Network intrusion detection is one of the most studied applications of machine
learning in cybersecurity. The core challenge is that in real-world networks,
attack traffic represents a tiny fraction of total traffic---often less than 1\%
---and many attacks, particularly zero-day exploits, have never been seen before.
This makes supervised classification insufficient as a sole defence: a model
that has only seen known attacks will fail to detect novel ones.

Unsupervised anomaly detection addresses this by learning only what \emph{normal}
traffic looks like, then flagging anything that deviates as potentially malicious.
The fundamental insight is that a model trained exclusively on normal data will
fail to reconstruct or generate anomalous samples well, and this failure itself
becomes the detection signal.

\subsection{Baseline Approaches and Their Limitations}

The two dominant unsupervised paradigms for network anomaly detection on KDD99
are reconstruction-based autoencoders \citep{ae_analysis_nids} and generative
adversarial networks \citep{alad}. Both exhibit structural weaknesses that
motivate this proposal:

\begin{enumerate}[leftmargin=2em]
    \item \textbf{Unstructured latent space (AE)}: Standard autoencoders impose
          no constraint on the bottleneck representation \citep{ae_analysis_nids}.
          A 4-dimensional latent code can take any shape; there is no principle
          by which a latent code can be declared ``normal-looking'' without
          re-decoding through the full decoder.
    \item \textbf{Training instability (GAN)}: GANs operating on the full
          114-dimensional input space are prone to convergence failure
          \citep{alad}. The discriminator is trained to separate
          \emph{generated} from \emph{real} samples---a proxy for normality,
          not a direct measure of it.
    \item \textbf{Uniform feature treatment (both)}: All 114 features receive
          equal weight during encoding. In KDD99, many one-hot encoded service
          and flag features are near-zero for most samples, while features such
          as \texttt{serror\_rate} and \texttt{dst\_host\_count} are highly
          discriminative for specific attack families. No existing approach on
          this dataset learns per-feature importance adaptively from the
          model's own reconstruction behaviour.
\end{enumerate}

\subsection{Proposed Contribution}

We propose a new architectural mechanism---\textbf{Two-Pass
Reconstruction-Error-Guided Attention (RE-Attention)}---embedded in an
Adversarial Autoencoder framework. A literature search (Section~\ref{sec:related})
confirms that while adversarial AEs, attention AEs, and dual scoring have all
been explored separately and in partial combinations, the specific mechanism of
using the per-feature reconstruction error vector as the attention signal for a
second encoding pass has not been reported in any prior work on network anomaly
detection.

% =============================================================================
\section{Dataset Description}
% =============================================================================

\subsection{KDD Cup 1999 (KDD99)}

The KDD Cup 1999 dataset~\citep{tavallaee2009detailed} was created from the
DARPA 1998 network intrusion dataset and contains approximately 5 million
network connection records. Each record describes a single TCP/IP connection
with 41 features: 34 continuous (traffic statistics such as byte counts, error
rates, and connection durations) and 7 categorical (protocol type, service,
connection flag, and binary indicators).

After one-hot encoding the categorical variables, the feature dimension expands
to \textbf{114 features}. The dataset contains 23 distinct attack types grouped
into four families:

\begin{table}[h!]
\centering
\caption{KDD99 Attack Families and Representative Types}
\begin{tabularx}{\textwidth}{lXl}
\toprule
\textbf{Family} & \textbf{Description} & \textbf{Examples} \\
\midrule
DoS   & Denial-of-service; overwhelm the target server
       & neptune, smurf, back, teardrop \\
Probe & Silent network scanning; reconnaissance before attack
       & ipsweep, nmap, satan, portsweep \\
R2L   & Remote-to-local; attempt unauthorised access from network
       & warezclient, guess\_passwd, ftp\_write \\
U2R   & User-to-root; attempt to escalate privileges to root
       & rootkit, buffer\_overflow, loadmodule \\
\bottomrule
\end{tabularx}
\end{table}

\noindent
Following standard experimental practice in this domain \citep{ae_analysis_nids,alad},
the dataset is subsampled to contain only 1\% anomalies relative to normal traffic,
simulating the realistic imbalance encountered in operational networks.
The preprocessing pipeline one-hot encodes categorical variables, applies
\texttt{MinMaxScaler} normalisation, and splits the data 75/25 into training
and test sets.

\subsection{NSL-KDD}

The NSL-KDD dataset~\citep{tavallaee2009detailed} addresses known deficiencies
in KDD99: duplicate records that cause classifiers to over-fit to common
patterns are removed, and the selection of records from each difficulty level
is inversely proportional to their prevalence in the original dataset, making
rare attack families harder to detect. It serves as the secondary evaluation
benchmark in this study, consistent with its use in prior work
\citep{vae_attention_lstm,ae_analysis_nids}.

% =============================================================================
\section{Related Work}
\label{sec:related}
% =============================================================================

The literature on deep learning for network anomaly detection is extensive.
We focus on the work most relevant to the proposed architecture, organised by
the primary technique used.

\subsection{Autoencoders for Network Intrusion Detection}

Autoencoders have been the dominant deep learning approach for unsupervised
network anomaly detection. \citet{ae_analysis_nids} provide a systematic
comparison of vanilla AE, denoising AE, sparse AE, and variational AE on the
KDD99 and NSL-KDD datasets, establishing that reconstruction error is a reliable
anomaly signal when training is restricted to normal traffic. \citet{improving_ae_nslkdd}
specifically improve standard AE performance on NSL-KDD through architectural
choices, providing an upper bound for vanilla AE baselines.

\textbf{Relevance to this work}: These papers establish the AE baseline (our
Condition~1) and confirm the reconstruction-error paradigm on the exact datasets
used in this study.

\subsection{Adversarial Regularisation of Autoencoders}

A key limitation of vanilla AEs is the absence of any constraint on the latent
space. Two papers address this for general anomaly detection:

\citet{alad} propose \textbf{Adversarially Learned Anomaly Detection (ALAD)},
a bi-directional GAN framework where an encoder maps data to latent codes
simultaneously with the generator mapping latent codes to data. Two discriminators
enforce cycle-consistency in both the data space and the latent space. ALAD uses
dual anomaly scores based on reconstruction error and discriminator outputs,
achieving state-of-the-art results on KDD99 and several image datasets with
454 citations as of 2026. ALAD is the closest prior work to the adversarial
regularisation component we propose.

\citet{arcade} present \textbf{ARCADE}, an adversarially regularised convolutional
AE specifically for network anomaly detection. ARCADE is trained exclusively on
normal traffic and uses an adversarial training strategy to reduce the AE's
ability to reconstruct out-of-distribution flows, evaluated on the CICIDS2017
dataset. Crucially, ARCADE uses no attention mechanism and no dual anomaly
scoring.

\citet{wavelet_aae} combine Haar wavelet transforms with an adversarial
autoencoder for detecting malicious TLS-encrypted network traffic, demonstrating
that adversarial AE consistently outperforms vanilla AE in the network security
domain.

\textbf{Gap}: ALAD, ARCADE, and the wavelet-AAE work all use input-driven or
data-space adversarial training. None employ attention guided by reconstruction
error.

\subsection{Attention Mechanisms for Intrusion Detection}

Attention mechanisms have been integrated with autoencoders and sequence models
for network intrusion detection, but almost exclusively as temporal attention
over LSTM hidden states.

\citet{vae_attention_lstm} propose a multi-model architecture combining a
class-wise focal-loss variational AE (\textbf{CWFL-VAE}) with an attention
mechanism-based LSTM (\textbf{AM-LSTM}) for imbalanced network traffic on
NSL-KDD and CSE-CIC-IDS2018. The attention is temporal (over LSTM time steps),
not feature-level, and the VAE is used for data balancing rather than direct
anomaly scoring.

\citet{attention_ae_latent} embed an attention module at the bottleneck of an
AE to selectively enhance important dimensions of the latent vector before
decoding. The attention is applied to the \emph{latent code}, not the input
features, and is not conditioned on reconstruction error.

\textbf{Gap}: Existing attention approaches for network intrusion are either
(a) temporal attention over sequence models, or (b) bottleneck attention on
the latent code. \textbf{No prior work uses the per-feature reconstruction
error vector as the attention signal for a second encoding pass.}

\subsection{Combined and Hybrid Approaches}

\citet{aladaen} propose the most comprehensive hybrid to date:
\textbf{ALADAEN}, an Active Learning-assisted Attention Adversarial Dual
AutoEncoder. The framework combines active learning (oracle querying for
uncertain samples), an adversarial dual autoencoder, and attention mechanisms
for Advanced Persistent Threat (APT) detection on DARPA Transparent Computing
provenance traces (not KDD99). While structurally ambitious, ALADAEN's active
learning loop requires an oracle and is evaluated on a different domain. The
attention mechanism is not specified as error-driven.

\citet{vae_wgan} combine a VAE and Wasserstein GAN into a
\textbf{VAE-WGAN} hybrid for network intrusion detection on KDD99. The
generative component is used primarily for data augmentation of minority
attack classes, and the anomaly detection uses a downstream LSTM-MSCNN
classifier rather than the latent representation directly.

\subsection{Summary of Gaps}

Table~\ref{tab:related_gap} summarises the key prior works and identifies the
architectural features missing from each that motivate this proposal.

\begin{table}[h!]
\centering
\caption{Comparison of Related Works and the Proposed Method}
\label{tab:related_gap}
\resizebox{\textwidth}{!}{%
\begin{tabular}{lccccc}
\toprule
\textbf{Method} & \textbf{Year} & \textbf{Adv. Latent}
  & \textbf{Attention} & \textbf{Dual Score} & \textbf{KDD99/NSL} \\
\midrule
Vanilla AE \citep{ae_analysis_nids}        & 2021 & \texttimes & \texttimes & \texttimes & \checkmark \\
ALAD \citep{alad}                          & 2018 & \checkmark & \texttimes & \checkmark & \checkmark \\
ARCADE \citep{arcade}                      & 2022 & \checkmark & \texttimes & \texttimes & \texttimes \\
Wavelet+AAE \citep{wavelet_aae}            & 2018 & \checkmark & \texttimes & \texttimes & \texttimes \\
VAE+Attn-LSTM \citep{vae_attention_lstm}   & 2025 & \texttimes & Temporal   & \texttimes & \checkmark \\
Attn-AE (bottleneck) \citep{attention_ae_latent} & 2021 & \texttimes & Latent & \texttimes & \texttimes \\
ALADAEN \citep{aladaen}                    & 2025 & \checkmark & Unspecified& \checkmark & \texttimes \\
VAE-WGAN \citep{vae_wgan}                  & 2024 & \checkmark & \texttimes & \texttimes & \checkmark \\
\midrule
\textbf{Proposed (RE-Attn AAE)}            & 2026 & \checkmark
  & \textcolor{novelred}{\textbf{Error-guided}} & \checkmark & \checkmark \\
\bottomrule
\end{tabular}}
\end{table}

% =============================================================================
\section{Problem Formulation}
% =============================================================================

Let $\mathbf{x} \in \mathbb{R}^{d}$ denote a network connection feature vector
($d = 114$ after one-hot encoding). The training set
$\mathcal{D}_{\text{train}} = \{\mathbf{x}_i\}_{i=1}^{N}$ contains predominantly
normal traffic (99\% normal, 1\% anomalous). The test set
$\mathcal{D}_{\text{test}}$ contains an unknown mixture of normal and
anomalous connections with ground-truth binary labels
$y_i \in \{0\ (\text{normal}),\ 1\ (\text{anomaly})\}$.

The objective is to learn a function
$f: \mathbb{R}^{d} \rightarrow \mathbb{R}_{\geq 0}$ that assigns an anomaly
score $s_i = f(\mathbf{x}_i)$ such that a threshold $\tau$ maximises
the area under the ROC curve (AUC-ROC) and the F1-score on
$\mathcal{D}_{\text{test}}$.

\noindent\textbf{Key Challenges}:
\begin{enumerate}[leftmargin=2em]
    \item \textit{Class imbalance}: anomalies comprise $\leq 1\%$ of traffic.
    \item \textit{Feature heterogeneity}: 114 features vary greatly in
          discriminative power; uniform weighting dilutes the signal from
          informative features.
    \item \textit{Latent space structure}: without a prior on the latent
          distribution, anomaly scoring in latent space is ill-defined.
    \item \textit{Multi-family attacks}: four attack families (DoS, Probe,
          R2L, U2R) have distinct signatures; a single threshold may fail on
          the rarer families (U2R, R2L).
\end{enumerate}

% =============================================================================
\section{Proposed Methodology}
% =============================================================================

\subsection{Overview}

We propose the \textbf{RE-Attention Adversarial Autoencoder (RE-Attn-AAE)},
an architecture built on three components that address the four challenges stated
above:

\begin{enumerate}[leftmargin=2em]
    \item \textbf{Two-Pass Reconstruction-Error-Guided Attention (RE-Attention)}:
          addresses feature heterogeneity by computing attention weights from
          where the model is currently failing, not from the raw input.
    \item \textbf{Adversarial Gaussian Latent Regularisation}: addresses latent
          space structure by enforcing $\mathbf{z} \approx \mathcal{N}(0, I)$
          through a small latent discriminator.
    \item \textbf{Dual Anomaly Scoring with learnable blend}: addresses class
          imbalance and multi-family attacks by combining two complementary
          signals: reconstruction error and latent normality.
\end{enumerate}

\subsection{Architecture}

\subsubsection{Pass 1 --- Standard Reconstruction}

The first pass is a standard deep autoencoder following the architecture
established for KDD99 in \citet{ae_analysis_nids}:

\begin{align}
\mathbf{z}^{(1)} &= \text{Enc}^{(1)}(\mathbf{x})
    \quad \in \mathbb{R}^{4} \label{eq:enc1}\\
\hat{\mathbf{x}} &= \text{Dec}(\mathbf{z}^{(1)})
    \quad \in \mathbb{R}^{114} \label{eq:dec}\\
\mathbf{e} &= (\mathbf{x} - \hat{\mathbf{x}})^2
    \quad \in \mathbb{R}^{114} \label{eq:error}
\end{align}

\noindent
where $\text{Enc}^{(1)}$ has layers
$114 \to 96 \to 64 \to 48 \to 16 \to 4$ (all tanh, with 10\% dropout),
and $\text{Dec}$ is its mirror image $4 \to 16 \to 48 \to 64 \to 96 \to 114$.

\subsubsection{RE-Attention Layer}

The per-feature squared reconstruction error $\mathbf{e}$ is passed through
a small two-layer network to produce a soft attention mask:

\begin{equation}
\mathbf{a} = \sigma\!\left(W_2 \cdot \text{ReLU}(W_1 \mathbf{e} + \mathbf{b}_1) + \mathbf{b}_2\right)
    \quad \in [0,1]^{114}
\label{eq:attention}
\end{equation}

\noindent
where $W_1 \in \mathbb{R}^{64 \times 114}$, $W_2 \in \mathbb{R}^{114 \times 64}$,
and $\sigma$ is the sigmoid function. The attended input is:

\begin{equation}
\tilde{\mathbf{x}} = \mathbf{x} \odot \mathbf{a}
\label{eq:attended}
\end{equation}

\noindent
\textbf{Interpretation}: A high reconstruction error on feature $j$ in Pass~1
produces a large $e_j$, which pushes $a_j$ toward 1, preserving that feature
in $\tilde{\mathbf{x}}$ for the second pass. Features that are already
well-reconstructed (low $e_j$) are attenuated. The model learns
\emph{where it is failing} and directs encoding capacity there.

\subsubsection{Pass 2 --- Adversarially Regularised Re-Encoding}

The attended input $\tilde{\mathbf{x}}$ is encoded by a second encoder
of the same architecture to produce the refined latent code:

\begin{equation}
\mathbf{z}^{(2)} = \text{Enc}^{(2)}(\tilde{\mathbf{x}})
    \quad \in \mathbb{R}^{4}
\label{eq:enc2}
\end{equation}

\noindent
A small latent discriminator $D_\text{lat}$ enforces a Gaussian prior on
$\mathbf{z}^{(2)}$:

\begin{equation}
D_\text{lat}(\mathbf{z}) = \sigma(W_3 \cdot \text{ReLU}(W_2' \cdot \text{ReLU}(W_1' \mathbf{z})))
    : \mathbb{R}^4 \to [0,1]
\label{eq:dlat}
\end{equation}

\noindent
where the layers are $4 \to 32 \to 16 \to 1$ (a small and stable network
given the 4-dimensional input).

\subsection{Training Procedure}

Training alternates between three phases per mini-batch:

\begin{tcolorbox}[
  title={\textbf{Algorithm 1:} RE-Attn-AAE Training (one mini-batch)},
  colback=lightgray, colframe=darkblue, fonttitle=\small\bfseries,
  breakable, enhanced
]
\small
\textbf{Input:} Batch $\{\mathbf{x}_i\}_{i=1}^{B}$ from $\mathcal{D}_\text{train}$

\medskip
\textbf{Phase 1 --- Reconstruction}
\begin{enumerate}[leftmargin=1.5em, topsep=2pt, itemsep=1pt]
  \item $\hat{\mathbf{x}} \leftarrow \text{Dec}(\text{Enc}^{(1)}(\mathbf{x}))$
  \item $\mathbf{e} \leftarrow (\mathbf{x} - \hat{\mathbf{x}})^2$
        \hfill \textit{// per-feature squared error}
  \item $\mathbf{a} \leftarrow \text{RE-Attn}(\mathbf{e})$;\quad
        $\tilde{\mathbf{x}} \leftarrow \mathbf{x} \odot \mathbf{a}$
  \item $\hat{\mathbf{x}}^{(2)} \leftarrow \text{Dec}(\text{Enc}^{(2)}(\tilde{\mathbf{x}}))$
  \item $\mathcal{L}_\text{recon} \leftarrow
        \|\mathbf{x} - \hat{\mathbf{x}}\|^2 +
        \|\mathbf{x} - \hat{\mathbf{x}}^{(2)}\|^2$
  \item Update $\text{Enc}^{(1)}, \text{Dec}, \text{RE-Attn}, \text{Enc}^{(2)}$
        via $\nabla \mathcal{L}_\text{recon}$
\end{enumerate}

\textbf{Phase 2 --- Latent Discriminator Update}
\begin{enumerate}[leftmargin=1.5em, topsep=2pt, itemsep=1pt]
  \item Sample $\mathbf{z}_\text{real} \sim \mathcal{N}(0, I_4)$
  \item $\mathbf{z}_\text{enc} \leftarrow \text{Enc}^{(2)}(\tilde{\mathbf{x}})$
        \hfill \textit{// encoder output (frozen here)}
  \item $\mathcal{L}_D \leftarrow
        -\mathbb{E}[\log D_\text{lat}(\mathbf{z}_\text{real})]
        -\mathbb{E}[\log(1 - D_\text{lat}(\mathbf{z}_\text{enc}))]$
  \item Update $D_\text{lat}$ via $\nabla \mathcal{L}_D$
\end{enumerate}

\textbf{Phase 3 --- Encoder Adversarial Update}
\begin{enumerate}[leftmargin=1.5em, topsep=2pt, itemsep=1pt]
  \item $\mathcal{L}_G \leftarrow
        -\mathbb{E}[\log D_\text{lat}(\text{Enc}^{(2)}(\tilde{\mathbf{x}}))]$
  \item Update $\text{Enc}^{(2)}$ via $\nabla \mathcal{L}_G$
        \hfill \textit{// encoder learns to fool $D_\text{lat}$}
\end{enumerate}
\end{tcolorbox}

\subsection{Anomaly Scoring}

At test time, for each sample $\mathbf{x}$, two scores are computed:

\begin{align}
s_\text{recon} &= \tfrac{1}{d}\|\mathbf{x} - \hat{\mathbf{x}}\|^2
    \label{eq:score_recon}\\
s_\text{latent} &= 1 - D_\text{lat}(\mathbf{z}^{(2)})
    \label{eq:score_latent}\\
s_\text{final}(\alpha) &= \alpha \cdot s_\text{recon}
    + (1 - \alpha) \cdot s_\text{latent},\quad \alpha \in [0, 1]
    \label{eq:score_final}
\end{align}

\noindent
The reconstruction score $s_\text{recon}$ captures attacks that distort many
features simultaneously (e.g., DoS). The latent score $s_\text{latent}$
captures subtle attacks whose latent code falls outside the Gaussian prior
even if individual features appear normal (e.g., R2L, U2R). The $\alpha$
parameter is swept to find the optimal blend, yielding an interpretable
curve of AUC-ROC vs.\ $\alpha$.

% =============================================================================
\section{Experimental Design}
% =============================================================================

\subsection{Ablation Conditions}

The experiment is a five-condition ablation study where each condition adds
exactly one architectural component. This isolates the contribution of each
component independently.

\begin{table}[h!]
\centering
\caption{Five Experimental Conditions and Their Architectural Components}
\label{tab:conditions}
\begin{tabularx}{\textwidth}{clXXXX}
\toprule
\textbf{C\#} & \textbf{Name} & \textbf{RE-Attn} & \textbf{Adv. Prior}
  & \textbf{Dual Score} & \textbf{Tests} \\
\midrule
C1 & Vanilla AE \citep{ae_analysis_nids}  & \texttimes & \texttimes & \texttimes
   & Reconstruction-only lower bound \\
C2 & GAN discriminator \citep{alad}       & \texttimes & \texttimes & \texttimes
   & Discriminator-only lower bound \\
C3 & AAE (no attention)   & \texttimes & \checkmark & Recon only
   & Does Gaussian prior improve reconstruction-based detection? \\
C4 & AAE + Dual Score     & \texttimes & \checkmark & \checkmark
   & Does the latent score add signal beyond reconstruction? \\
C5 & RE-Attn-AAE (full)   & \checkmark & \checkmark & \checkmark
   & Does error-guided attention further improve C4? \\
\bottomrule
\end{tabularx}
\end{table}

\noindent
The difference C3 $\to$ C4 tests the value of the dual scoring mechanism.
The difference C4 $\to$ C5 tests the specific novel contribution:
\textbf{error-guided attention}.

\subsection{Evaluation Metrics}

For each condition the following metrics are computed on the test set:

\begin{itemize}[leftmargin=2em]
    \item \textbf{AUC-ROC}: area under the receiver operating characteristic
          curve; measures overall ranking quality.
    \item \textbf{AUC-PR}: area under the precision-recall curve; more
          informative for the highly imbalanced 1\% anomaly setting.
    \item \textbf{F1-score}: at the threshold that maximises F1 on the test set.
    \item \textbf{False Negative Rate (FNR)}: proportion of attacks missed;
          critical for security applications.
    \item \textbf{Per-family AUC}: AUC-ROC computed separately for DoS, Probe,
          R2L, and U2R to reveal which attack families benefit most from
          error-guided attention.
\end{itemize}

\subsection{Additional Analysis: $\alpha$-Sweep}

For conditions C4 and C5, the blend parameter $\alpha$ is swept from 0 to 1 in
steps of 0.05. A plot of AUC-ROC vs.\ $\alpha$ for each condition directly
reveals:
\begin{itemize}[leftmargin=2em]
    \item The optimal $\alpha$ for each condition.
    \item Whether the dual score is additive (peak strictly between 0 and 1)
          or dominated by one signal.
    \item Whether RE-Attention changes the optimal blend point, indicating that
          it creates a qualitatively different latent representation.
\end{itemize}

\subsection{Interpretability Analysis: Attention Weight Visualisation}

After training C5, the attention weights $\mathbf{a}$ are recorded for
\emph{all correctly detected test anomalies} grouped by attack family. The
mean attention weight per feature per family is plotted as a heatmap
(114 features $\times$ 4 attack families). This answers the question:
\emph{Which network traffic features does the model focus on when it detects
each attack type?}

This is the key interpretability result. Existing attention works
\citep{vae_attention_lstm,attention_ae_latent} do not provide this feature-level
analysis on KDD99/NSL-KDD attack families.

\subsection{Hyperparameters}

\begin{table}[h!]
\centering
\caption{Shared Hyperparameters Across All Conditions}
\begin{tabular}{ll}
\toprule
\textbf{Parameter} & \textbf{Value} \\
\midrule
Batch size & 512 \\
Optimiser & Adam ($\beta_1 = 0.5$, $\beta_2 = 0.999$) \\
Learning rate & $1 \times 10^{-5}$ \\
Epochs & 30 (C1--C5) \\
Latent dimension ($d_z$) & 4 \\
RE-Attention hidden size & 64 \\
Latent discriminator & $4 \to 32 \to 16 \to 1$ \\
Anomaly \% in training & 1\% \\
Random seed & 42 \\
\bottomrule
\end{tabular}
\end{table}

\subsection{Infrastructure}

All experiments run in a single Kaggle notebook on a T4 GPU
(\texttt{tensorflow==2.x}, \texttt{scikit-learn}, \texttt{numpy}).
Estimated total runtime: $\approx$35 minutes for all five conditions at 30 epochs,
well within Kaggle's 9-hour session limit and the 2-hour target.

% =============================================================================
\section{Expected Results and Hypotheses}
% =============================================================================

\begin{enumerate}[leftmargin=2em, label=\textbf{H\arabic*.}]

    \item \textbf{Adversarial prior improves over vanilla AE (C3 vs C1)}:
          Enforcing a Gaussian prior on the latent space should improve AUC-PR
          because anomalous latent codes are more likely to fall outside the
          regularised distribution, even before adding the explicit latent score.
          We expect C3 AUC-PR $>$ C1 AUC-PR.

    \item \textbf{Dual scoring adds complementary signal (C4 vs C3)}:
          The reconstruction score alone is most effective for DoS attacks
          (which distort many features). The latent score should add detection
          power for Probe and R2L attacks (which are subtler in feature space).
          We expect a distinct peak in the AUC-ROC vs.\ $\alpha$ curve at
          $0 < \alpha^* < 1$ for C4, confirming that both signals contribute.

    \item \textbf{RE-Attention further improves over input-driven alternatives (C5 vs C4)}:
          Error-guided attention directs the second-pass encoder to concentrate
          on the most anomaly-relevant features. We expect this to improve
          per-family AUC for the rarer and harder attack families (R2L, U2R),
          where the anomaly signature is concentrated in a small number of
          features that the error vector will highlight.

    \item \textbf{Attention weights are interpretable and consistent with domain
          knowledge}: We expect the heatmap to show high attention on
          \texttt{serror\_rate}, \texttt{count}, \texttt{src\_bytes} for DoS;
          on \texttt{dst\_host\_count}, \texttt{dst\_host\_srv\_count} for Probe;
          on \texttt{num\_failed\_logins}, \texttt{is\_guest\_login} for R2L;
          and on \texttt{root\_shell}, \texttt{num\_root} for U2R. Agreement
          with known network attack signatures validates the attention mechanism.

\end{enumerate}

% =============================================================================
\section{Limitations and Risks}
% =============================================================================

\begin{itemize}[leftmargin=2em]
    \item \textbf{KDD99 is a synthetic dataset}: It was generated from a
          simulated military network in 1998. Results may not transfer directly
          to modern network traffic. NSL-KDD reduces but does not eliminate
          this concern.

    \item \textbf{Latent dimension}: The 4-dimensional bottleneck is very
          small for the dual encoder design. If the second encoder does not
          converge stably, the latent discriminator may receive noisy gradients.
          A larger latent dimension (8 or 16) is a contingency.

    \item \textbf{Three-phase training instability}: Alternating three
          gradient updates per batch is standard for adversarial AEs but
          can be sensitive to the relative learning rates of encoder and
          discriminator. Careful monitoring of discriminator loss is required.

    \item \textbf{Dual encoder parameter count}: The full RE-Attn-AAE has
          approximately double the encoder parameters of the vanilla AE baseline
          (C1). This is controlled for in the ablation by comparing all conditions
          at the same number of training epochs.
\end{itemize}

% =============================================================================
\section{Conclusion}
% =============================================================================

This proposal describes a focused architectural experiment for network anomaly
detection on the KDD99 and NSL-KDD datasets. The proposed \textbf{RE-Attention
Adversarial Autoencoder} advances beyond established reconstruction-based
\citep{ae_analysis_nids} and adversarial \citep{alad,arcade} baselines through
three principled contributions: adversarial Gaussian regularisation of the latent
space, dual anomaly scoring combining reconstruction and latent normality signals,
and a novel error-guided attention mechanism that uses the first-pass reconstruction
failure to re-weight the input for a more targeted second-pass encoding.

A literature search confirmed that while adversarial AEs, attention AEs, and
dual scoring have been explored individually and in partial combinations
\citep{alad,arcade,aladaen,wavelet_aae,vae_attention_lstm,attention_ae_latent},
the specific mechanism of computing attention weights from the per-feature
reconstruction error vector has not been reported for network intrusion
detection. This constitutes the novel contribution of the experiment.

The controlled five-condition ablation study allows each component's contribution
to be measured independently. The interpretability analysis --- mapping attention
weights to KDD99 feature names and attack families --- provides a qualitative
result that complements the quantitative metrics and may offer genuine insight
into which aspects of network traffic most strongly signal each attack type.

% =============================================================================
% References
% =============================================================================
\bibliographystyle{plainnat}

\begin{thebibliography}{99}

\bibitem[Zenati et al.(2018)]{alad}
Zenati, H., Romain, M., Foo, C.-S., Lecouat, B., \& Chandrasekhar, V. (2018).
\textit{Adversarially Learned Anomaly Detection}.
IEEE International Conference on Data Mining (ICDM).
\url{https://arxiv.org/abs/1812.02288}

\bibitem[Lunardi et al.(2022)]{arcade}
Lunardi, W.\ T., Lopez, M.\ A., \& Giacalone, J.-P. (2022).
\textit{ARCADE: Adversarially Regularized Convolutional Autoencoder for Network
Anomaly Detection}.
arXiv preprint arXiv:2205.01432.
\url{https://arxiv.org/pdf/2205.01432}

\bibitem[Benabderrahmane et al.(2025)]{aladaen}
Benabderrahmane, S., Cheney, J., \& Rahwan, T. (2025).
\textit{Ranking-enhanced anomaly detection using Active Learning-assisted
Attention Adversarial Dual AutoEncoder}.
\textit{Scientific Reports}, 15.
\url{https://doi.org/10.1038/s41598-025-25621-0}

\bibitem[Puuska et al.(2018)]{wavelet_aae}
Puuska, S., Kokkonen, T., Alatalo, J., \& Heilimo, E. (2018).
\textit{Anomaly-Based Network Intrusion Detection Using Wavelets and
Adversarial Autoencoders}.
International Conference on Security for Information Technology and
Communications, Springer.
\url{https://link.springer.com/chapter/10.1007/978-3-030-12942-2_18}

\bibitem[Abdulganiyu et al.(2025)]{vae_attention_lstm}
Abdulganiyu, O.\ H., Ait Tchakoucht, T., Saheed, Y., El Mouhtadi, M., \&
Alaoui, A. (2025).
\textit{Modified Variational Autoencoder and Attention Mechanism-Based Long
Short-Term Memory for Detecting Intrusions in Imbalanced Network Traffic}.
\textit{Security and Privacy}.
\url{https://doi.org/10.1002/spy2.70044}

\bibitem[Al-Qatf et al.(2021)]{attention_ae_latent}
Al-Qatf, M., et al. (2021).
\textit{Attention Autoencoder for Generative Latent Representational Learning
in Anomaly Detection}.
\textit{Sensors}, 22(1), 240.
\url{https://doi.org/10.3390/s22010240}

\bibitem[Zavrak \& Iskefiyeli(2021)]{ae_analysis_nids}
Zavrak, S., \& Iskefiyeli, M. (2021).
\textit{Analysis of Autoencoders for Network Intrusion Detection}.
\textit{Sensors}, 21(12), 4153.
\url{https://doi.org/10.3390/s21124153}

\bibitem[Tavallaee et al.(2009)]{tavallaee2009detailed}
Tavallaee, M., Bagheri, E., Lu, W., \& Ghorbani, A. A. (2009).
\textit{A Detailed Analysis of the KDD Cup 99 Data Set}.
IEEE CISDA.

\bibitem[Xu et al.(2024)]{vae_wgan}
Xu, Y., et al. (2024).
\textit{An intrusion detection method combining variational auto-encoder and
generative adversarial networks}.
\textit{Computer Networks}.
\url{https://doi.org/10.1016/j.comnet.2024.110728}

\bibitem[Hassan et al.(2025)]{arae_network}
Hassan, M., et al. (2025).
\textit{ARAE: An Adaptive Robust AutoEncoder for Network Anomaly Detection}.
\textit{Journal of Cyber Security}, 7(1).
\url{https://www.techscience.com/JCS/v7n1/65064/html}

\end{thebibliography}

\end{document}
